\section{Termín 30. ledna 2025}

\subsection{Rozvin funkce}
Nalezněte rozvoj funkce
\[
    f(x) = \frac{x^3+2x}{x^2 - 2}
\]
v okolí bodu $x_0 = 0$. Určete největší otevřený interval, kde nastává rovnost funkce a příslušné řady.

\subsection{Tečná rovina}
Nalezněte bod na ploše
\[
    z = 4 - x^2 - y^2 \text{, }
\]
ve kterém je tečná rovina rovnoběžná s rovinou $2x + 2y + z=0$. Určete tuto tečnou rovinu.

\subsection{Dvojný integrál}
Přepište následující integrál
\[
    \int_{0}^{1} \int_{-1}^{-\sqrt{2y - y^2}} f(x, y) \dif x \dif y
\]
\begin{enumerate}[a), noitemsep]
    \item v opačném pořadí integrace,
    \item v polárních souřadnicích se středem v počátku v pořadí $\dif \varrho \dif \varphi$.
\end{enumerate}


\subsection{Konzervativní vektorové pole}
Ověřte, že vektorové pole
\[
    \vec{F} (x, y, z) = (y \sin z - y^2 \sin x, 2y \cos x + x \sin z, \frac{1}{z} + xy \cos z)
\]
je konzervativní a nalezněte jeho potenciál. Určete práci pole z bodu $(1, 0, 1)$ do bodu $(0, 1, 2)$.

\subsection{Objem tělesa}
Určete objem tělesa vymezeného válcem $x^2 + y^2 = 1$ plochou $y = |z|$.